% Unit 037 terjemahan Bahasa Indonesia dari chapter3.tex baris 946--1060.
% Sumber: Methods of Algebra, Volume 2, commit
% 9a5803ff77dd3257484cb177f851a73770a59dd3 (CC BY 4.0).
% stable-unit-id: o014.aljabr2.chapter3.abelian-category-complexes
% Cakupan: Bagian 3.6 lengkap.

% segment-id: o014.li2.chapter3.abelian-category-complexes.h001
\section{Kompleks pada kategori abelian}\label{sec:Abel-cplx}
\hypertarget{o014.aljabr2.chapter3.abelian-category-complexes}{ }

% segment-id: o014.li2.chapter3.abelian-category-complexes.p001
Dalam bagian ini, $\mathcal{A}$ diasumsikan sebagai kategori abelian. Seperti
biasa, semua pernyataan tentang keaditifan dalam bagian ini juga dapat
ditingkatkan ke kasus ketika $\mathcal{A}$ merupakan kategori $\Bbbk$-linear.

% segment-id: o014.li2.chapter3.abelian-category-complexes.p002
Untuk suatu kompleks $X$ pada kategori abelian, kita dapat membahas
kohomologinya (Definisi~\ref{def:cohomology}). Ingat bahwa $[n]$ bukan hanya
menggeser indeks atas kompleks, melainkan juga mengalikan $d_X^\bullet$ dengan
tanda $(-1)^n$. Akan tetapi, perubahan tanda ini tidak mengubah kernel maupun
citra setiap $d_X$. Dengan demikian,
% segment-id: o014.li2.chapter3.abelian-category-complexes.q001
\[ \Hm^k\left( X[n] \right) = \Hm^{n+k} \left(X \right), \quad k, n \in \Z . \]

% segment-id: o014.li2.chapter3.abelian-category-complexes.p003
Proposisi~\ref{prop:additive-cat-cplx} telah memastikan bahwa
$\cate{C}(\mathcal{A})$ merupakan kategori aditif. Berikut ini ditunjukkan bahwa
kategori tersebut juga abelian.

% segment-id: o014.li2.chapter3.abelian-category-complexes.pr001
\begin{proposisi}\label{prop:abelian-cat-cplx}
	Kategori $\cate{C}(\mathcal{A})$ merupakan kategori abelian. Secara lebih
	tepat, untuk semua $X,Y\in\Obj(\cate{C}(\mathcal{A}))$, berlaku
	$X \oplus Y = (X^n \oplus Y^n, (d_X^n, d_Y^n))_{n \in \Z}$, dan untuk
	setiap morfisme $f:X\to Y$ dapat diambil
% segment-id: o014.li2.chapter3.abelian-category-complexes.q002
	\begin{gather*}
		\Ker(f) = \left( \Ker(f^n) \right)_{n \in \Z}, \quad
		\Coker(f) = \left( \Coker(f^n) \right)_{n \in \Z}, \\
		\Image(f) = \left( \Image(f^n) \right)_{n \in \Z}, \quad
		\Coim(f) = \left( \Coim(f^n) \right)_{n \in \Z}.
	\end{gather*}%
	\footnote{Catatan penerjemah (O014-C040): sumber menuliskan
	$\Ker(f)^n$, $\Coker(f)^n$, $\Image(f)^n$, dan $\Coim(f)^n$ pada ruas
	kiri, tetapi menyamakannya dengan seluruh keluarga berindeks $n$ pada ruas
	kanan. Target menghapus pangkat $n$ pada ruas kiri agar keempatnya merupakan
	identitas kompleks; komponen derajat $n$ tetap tepat seperti yang
	ditampilkan pada ruas kanan.}

	Misalkan morfisme $f$ dan $g$ dalam $\cate{C}(\mathcal{A})$ dapat
	dikomposisikan dan $gf=0$. Maka
	$H:=\Hm\left[X\xrightarrow{f}Y\xrightarrow{g}Z\right]$ adalah kompleks
	berikut:
% segment-id: o014.li2.chapter3.abelian-category-complexes.l001
	\begin{itemize}
% segment-id: o014.li2.chapter3.abelian-category-complexes.i001
		\item suku derajat $n$ ialah
		$H^n:=\Hm\left[X^n\xrightarrow{f^n}Y^n\xrightarrow{g^n}Z^n\right]$;
% segment-id: o014.li2.chapter3.abelian-category-complexes.i002
		\item morfisme $d_H^n:H^n\to H^{n+1}$ ditentukan oleh $d_X^n$,
		$d_Y^n$, dan $d_Z^n$, bersama dengan funktorialitas $\Hm[\cdots]$
		(Proposisi~\ref{prop:homology-functoriality}).\footnote{Catatan
		penerjemah (O014-C041): sumber hanya menyebut $d_X^n$ dan $d_Y^n$.
		Target menambahkan $d_Z^n$, yang diperlukan untuk memberikan morfisme
		antara kedua diagram tiga-suku dan dengan demikian menginduksi $d_H^n$.}
	\end{itemize}
	Sebagai akibatnya, $X\xrightarrow{f}Y\xrightarrow{g}Z$ eksak jika dan
	hanya jika $X^n\xrightarrow{f^n}Y^n\xrightarrow{g^n}Z^n$ eksak untuk
	setiap $n\in\Z$.
\end{proposisi}

% segment-id: o014.li2.chapter3.abelian-category-complexes.pf001
\begin{bukti}
	Proposisi~\ref{prop:additive-cat-cplx} telah menjelaskan cara menggunakan
	funktor pelupa dari kompleks ke objek bergradasi,
	$U:\cate{C}(\mathcal{A})\to\mathcal{A}^{\Z}$, untuk mereduksi
	$\varinjlim$ dan $\varprojlim$ yang diperlukan menjadi konstruksi derajat
	demi derajat di $\mathcal{A}$. Pokoknya ialah bahwa $U$ menciptakan
	$\varinjlim$ dan $\varprojlim$ (Lema~\ref{prop:limits-cplx}). Hal ini
	memberikan konstruksi derajat demi derajat bagi $X\oplus Y$, $\Ker(f)$,
	$\Coker(f)$, dan seterusnya.

	Secara khusus, berdasarkan karakterisasi~\eqref{eqn:strict-morphism},
	morfisme kanonik $\Coim(f)\to\Image(f)$ diberikan oleh morfisme
	$\Coim(f^n)\to\Image(f^n)$ di $\mathcal{A}$. Isomorfisme derajat demi
	derajat adalah isomorfisme kompleks. Jadi, setiap morfisme dalam
	$\cate{C}(\mathcal{A})$ bersifat ketat. Oleh karena itu,
	$\cate{C}(\mathcal{A})$ merupakan kategori abelian.\footnote{Catatan
	penerjemah (O014-C042): kalimat penutup sumber menyebut $\mathcal{A}$,
	padahal kategori itu telah diasumsikan abelian, sedangkan kategori yang baru
	saja dibuktikan abelian ialah $\cate{C}(\mathcal{A})$. Target memperbaiki nama kategori
	tersebut.} Deskripsi $\Hm[X\to Y\to Z]$ diperoleh dengan cara yang sama.
\end{bukti}

% segment-id: o014.li2.chapter3.abelian-category-complexes.r001
\begin{catatan}
	Dengan cara yang sama dapat dibuktikan bahwa kategori bikompleks
	$\cate{C}^2(\mathcal{A})$, bahkan kategori kompleks lipat-$k$
	$\cate{C}^k(\mathcal{A})$, tetap merupakan kategori abelian untuk
	$k\in\Z_{\geq 1}$; lihat Catatan~\ref{rem:k-fold-cplx}.
\end{catatan}

% segment-id: o014.li2.chapter3.abelian-category-complexes.p004
Selanjutnya, perhatikan suatu morfisme $f:X\to Y$ dalam
$\cate{C}(\mathcal{A})$. Untuk setiap $n\in\Z$,
Proposisi~\ref{prop:homology-functoriality} menentukan secara unik
$\Hm^n(f):\Hm^n(X)\to\Hm^n(Y)$ sedemikian sehingga diagram berikut komutatif:
% segment-id: o014.li2.chapter3.abelian-category-complexes.g001
\begin{equation}\label{eqn:induced-morphism-cohomology}\begin{tikzcd}
	\Ker(d_X^n) \arrow[twoheadrightarrow, r] \arrow[d, "\text{diinduksi oleh $f$}"'] & \Hm^n(X) \arrow[hookrightarrow, r] \arrow[d, "{\Hm^n(f)}"] & \Coker(d_X^{n-1}) \arrow[d, "\text{diinduksi oleh $f$}"] \\
	\Ker(d_Y^n) \arrow[twoheadrightarrow, r] & \Hm^n(Y) \arrow[hookrightarrow, r] & \Coker(d_Y^{n-1}) .
\end{tikzcd}\end{equation}

% segment-id: o014.li2.chapter3.abelian-category-complexes.p005
Untuk morfisme yang dapat dikomposisikan
$X\xrightarrow{f}Y\xrightarrow{g}Z$, karakterisasi ini langsung memberikan
$\Hm^n(gf)=\Hm^n(g)\Hm^n(f)$; selain itu,
$\Hm^n(\identity_X)=\identity_{\Hm^n(X)}$.

% segment-id: o014.li2.chapter3.abelian-category-complexes.pr002
\begin{proposisi}[Kohomologi sebagai funktor]\label{prop:H-functor}
	\index[sym1]{Hn}
	Untuk setiap $n\in\Z$, definisi di atas memberikan funktor aditif
	$\Hm^n:\cate{C}(\mathcal{A})\to\mathcal{A}$.
\end{proposisi}

% segment-id: o014.li2.chapter3.abelian-category-complexes.pf002
\begin{bukti}
	Sifat $\Hm^n(gf)=\Hm^n(g)\Hm^n(f)$ dan
	$\Hm^n(\identity)=\identity$ langsung diperoleh dari
	karakterisasi~\eqref{eqn:induced-morphism-cohomology}. Dengan cara yang
	sama diperoleh $\Hm^n(f_1+f_2)=\Hm^n(f_1)+\Hm^n(f_2)$, serta
	$\Hm^n(tf)=t\Hm^n(f)$ apabila $\mathcal{A}$ merupakan kategori
	$\Bbbk$-linear dan $t\in\Bbbk$.
\end{bukti}

% segment-id: o014.li2.chapter3.abelian-category-complexes.p006
Barisan eksak pendek antar-kompleks secara otomatis menginduksi barisan eksak
panjang pada kohomologi. Inilah bentuk paling dasar dari barisan eksak panjang
dalam teori homologi.

% segment-id: o014.li2.chapter3.abelian-category-complexes.pr003
\begin{proposisi}[Barisan eksak pendek menginduksi barisan eksak panjang]\label{prop:long-exact-sequence-ses}
	\index{barisan eksak panjang (long exact sequence)}
	Misalkan $0\to X\xrightarrow{f}Y\xrightarrow{g}Z\to 0$ merupakan barisan
	eksak pendek dalam $\cate{C}(\mathcal{A})$. Maka Teorema~\ref{prop:snake-lemma}
	memberikan barisan eksak kanonik berikut di $\mathcal{A}$:

% segment-id: o014.li2.chapter3.abelian-category-complexes.g002
	\begin{equation*}\begin{tikzcd}[row sep=large]
	\cdots \arrow[r] & \Hm^{n-1}(Y) \arrow[r, "{\Hm^{n-1}(g)}" inner sep=0.7em] \arrow[phantom, d, ""{coordinate, name=A}] & \Hm^{n-1}(Z) \arrow[dll, rounded corners, "{\delta^{n-1}}" description, to path={
		-- ([xshift=2ex]\tikztostart.east)
		|- (A) [near end]\tikztonodes
		-| ([xshift=-2ex]\tikztotarget.west)
		-- (\tikztotarget)}] \\
	\Hm^n(X) \arrow[r, "{\Hm^n(f)}"] & \Hm^n(Y) \arrow[r, "{\Hm^n(g)}"] \arrow[phantom, d, ""{coordinate, name=B}] & \Hm^n(Z) \arrow[dll, rounded corners, "{\delta^n}" description, to path={
		-- ([xshift=2ex]\tikztostart.east)
		|- (B) [near end]\tikztonodes
		-| ([xshift=-2ex]\tikztotarget.west)
		-- (\tikztotarget)}] \\
	\Hm^{n+1}(X) \arrow[r, "{\Hm^{n+1}(f)}"' inner sep=0.7em] & \Hm^{n+1}(Y) \arrow[r] & \cdots
	\end{tikzcd}\end{equation*}

	Barisan ini mempunyai funktorialitas berikut. Jika terdapat diagram
	komutatif dalam $\cate{C}(\mathcal{A})$ yang baris-barisnya eksak,
% segment-id: o014.li2.chapter3.abelian-category-complexes.g003
	\[\begin{tikzcd}[row sep=small]
		0 \arrow[r] & X \arrow[r] \arrow[d] & Y \arrow[r] \arrow[d] & Z \arrow[d] \arrow[r] & 0 \\
		0 \arrow[r] & \underline{X} \arrow[r] & \underline{Y} \arrow[r] & \underline{Z} \arrow[r] & 0
	\end{tikzcd}\]
	maka diperoleh diagram komutatif
% segment-id: o014.li2.chapter3.abelian-category-complexes.g004
	\[\begin{tikzcd}[column sep=small, row sep=small]
		\cdots \arrow[r] & \Hm^n(X) \arrow[d] \arrow[r] & \Hm^n(Y) \arrow[r] \arrow[d] & \Hm^n(Z) \arrow[r] \arrow[d] & \Hm^{n+1}(X) \arrow[r] \arrow[d] & \cdots \\
		\cdots \arrow[r] & \Hm^n(\underline{X}) \arrow[r] & \Hm^n(\underline{Y}) \arrow[r] & \Hm^n(\underline{Z}) \arrow[r] & \Hm^{n+1}(\underline{X}) \arrow[r] & \cdots .
	\end{tikzcd}\]
\end{proposisi}

% segment-id: o014.li2.chapter3.abelian-category-complexes.pf003
\begin{bukti}
	Untuk setiap $n\in\Z$ terdapat barisan eksak
	$\Coker(d_X^{n-1})\to\Coker(d_Y^{n-1})\to\Coker(d_Z^{n-1})\to0$.
	Memang, ambil bagian $\Coker$ dari barisan eksak dalam
	Teorema~\ref{prop:snake-lemma} yang diterapkan pada diagram
% segment-id: o014.li2.chapter3.abelian-category-complexes.g005
	\[\begin{tikzcd}
		0 \arrow[r] & X^{n-1} \arrow[r] \arrow[d] & Y^{n-1} \arrow[r] \arrow[d] & Z^{n-1} \arrow[r] \arrow[d] & 0 \\
		0 \arrow[r] & X^n \arrow[r] & Y^n \arrow[r] & Z^n \arrow[r] & 0
	\end{tikzcd}\]%
	\footnote{Catatan penerjemah (O014-C043): sumber menyatakan barisan
	$\Coker(d_X^n)\to\Coker(d_Y^n)\to\Coker(d_Z^n)\to0$, sedangkan panah
	vertikal pada diagram yang langsung menyusul ialah $d_X^{n-1}$,
	$d_Y^{n-1}$, dan $d_Z^{n-1}$. Target menyesuaikan indeks kokernel dengan
	diagram; pergantian nama indeks kemudian memberikan baris berindeks $n-2$
	yang dipakai di bawah.}
	Dengan cara yang sama juga terdapat barisan eksak
	$0\to\Ker(d_X^n)\to\Ker(d_Y^n)\to\Ker(d_Z^n)$. Kedua barisan itu dapat
	ditempatkan dalam diagram komutatif dengan baris-baris eksak berikut:
% segment-id: o014.li2.chapter3.abelian-category-complexes.g006
	\begin{equation}\label{eqn:long-exact-sequence-ses-aux}\begin{tikzcd}
		& \Coker(d_X^{n-2}) \arrow[r] \arrow[d] & \Coker(d_Y^{n-2}) \arrow[r] \arrow[d] & \Coker(d_Z^{n-2}) \arrow[d] \arrow[r] & 0 \\
		0 \arrow[r] & \Ker(d_X^n) \arrow[r] & \Ker(d_Y^n) \arrow[r] & \Ker(d_Z^n) &
	\end{tikzcd}\end{equation}
	Panah-panah vertikal masing-masing diinduksi oleh $d_X^{n-1}$,
	$d_Y^{n-1}$, dan $d_Z^{n-1}$, serta mempunyai faktorisasi epi--mono
	berbentuk
	$\Coker(d^{n-2})\stackrel{d^{n-1}}{\twoheadrightarrow}
	\Image(d^{n-1})\hookrightarrow\Ker(d^n)$ (subskrip dihilangkan).

	Menurut Lema~\ref{prop:mono-epi-ker-coker} (iii), epimorfisme
	$\Coker(d^{n-2})\twoheadrightarrow\Image(d^{n-1})$ menentukan kernel
	panah-panah vertikal dalam~\eqref{eqn:long-exact-sequence-ses-aux}; lalu
	\eqref{eqn:homology-3} berturut-turut memberikan $\Hm^{n-1}(X)$,
	$\Hm^{n-1}(Y)$, dan $\Hm^{n-1}(Z)$. Demikian pula, monomorfisme
	$\Image(d^{n-1})\hookrightarrow\Ker(d^n)$ menentukan kokernel
	panah-panah vertikal tersebut, yang berturut-turut memberikan $\Hm^n(X)$,
	$\Hm^n(Y)$, dan $\Hm^n(Z)$. Penerapan Teorema~\ref{prop:snake-lemma}
	sekarang menghasilkan morfisme penghubung
	$\delta^{n-1}:\Hm^{n-1}(Z)\to\Hm^n(X)$ beserta bagian barisan eksak
	panjang yang memuat suku-suku $\Hm^n$ dan $\Hm^{n-1}$. Funktorialitasnya
	diperoleh dengan menerapkan Catatan~\ref{rem:connecting-canonical}.
\end{bukti}

% segment-id: o014.li2.chapter3.abelian-category-complexes.d001
\begin{definisi}[Kuasi-isomorfisme]\label{def:quasi-isomorphism}
	\index{kuasi-isomorfisme (quasi-isomorphism)}
	Misalkan $f:X\to Y$ merupakan morfisme dalam $\cate{C}(\mathcal{A})$.
	Jika $\Hm^n(f)$ merupakan isomorfisme untuk setiap $n\in\Z$, maka $f$
	disebut \emph{kuasi-isomorfisme}.
\end{definisi}

% segment-id: o014.li2.chapter3.abelian-category-complexes.pr004
\begin{proposisi}\label{prop:null-homotopic-cohomology}
	Misalkan $f:X\to Y$ merupakan morfisme dalam $\cate{C}(\mathcal{A})$ dan
	$n\in\Z$. Jika $f$ homotopik nol, maka $\Hm^n(f)=0$.
\end{proposisi}

% segment-id: o014.li2.chapter3.abelian-category-complexes.pf004
\begin{bukti}
	Ambil $h\in\Hom^{-1}(X,Y)$ sedemikian sehingga $f=d_Yh+hd_X$.
	Komposisi
	$\Ker(d_X^n)\hookrightarrow X^n\xrightarrow{h^{n+1}d_X^n}Y^n$
	ialah nol. Di sisi lain, $d_Y^{n-1}h^n$ berfaktor melalui
	$\Image(d_Y^{n-1})$. Kedua fakta ini bersama-sama memberikan
	$\Hm^n(f)=0$.
\end{bukti}

% segment-id: o014.li2.chapter3.abelian-category-complexes.co001
\begin{korolari}
	Untuk setiap $n\in\Z$, funktor kohomologi
	$\Hm^n:\cate{C}(\mathcal{A})\to\mathcal{A}$ berfaktor secara unik melalui
	$\cate{K}(\mathcal{A})$. Oleh karena itu, konsep kuasi-isomorfisme dapat
	diperluas ke morfisme dalam $\cate{K}(\mathcal{A})$. Jika $f:X\to Y$
	merupakan isomorfisme dalam $\cate{K}(\mathcal{A})$, maka $f$ merupakan
	kuasi-isomorfisme.
\end{korolari}

% segment-id: o014.li2.chapter3.abelian-category-complexes.pf005
\begin{bukti}
	Gabungkan Proposisi~\ref{prop:homotopy-category-universal-property} dan
	Proposisi~\ref{prop:null-homotopic-cohomology}.
\end{bukti}
